\subsection{Gestión y trazabilidad}\label{gestiuxf3n-y-trazabilidad}

\subsubsection{Trazabilidad end-to-end e integración del ERS}\label{trazabilidad-end-to-end}

La trazabilidad end-to-end del SGCV-IA conecta cada requisito con su
origen (Fuente/Stakeholder) y con todo lo que se deriva de él (Caso de
Uso, Clase de dominio, Estado, Criterio BDD, Caso de Prueba conceptual),
conforme a la cadena mínima exigida por la guía de la PE5:
Fuente/Stakeholder $\rightarrow$ RF $\rightarrow$ CU $\rightarrow$ Clase
UML $\rightarrow$ Estado $\rightarrow$ Criterio BDD $\rightarrow$ Caso
de Prueba.

\paragraph{Verificación de correcciones de la auditoría anterior}

Se verificaron los dos hallazgos críticos de la ronda anterior sobre
\texttt{main.tex}, ambos confirmados como resueltos:

\begin{itemize}
\item \textbf{Colisión de ID (RF-25):} el identificador RF-25 quedó
  correctamente reservado para el requisito aprobado por el CCB en
  PE4 (confirmación de asistencia a citas), y RF-26 (Portal de
  ejercicio de derechos ARCO) fue reasignado sin conflicto.
\item \textbf{Contenido de otro dominio:} el texto ajeno al proyecto
  ("predicción de cosecha") detectado en la fila Could-have de la
  matriz MoSCoW fue eliminado; la fila describe correctamente
  "Ninguno (0 RF)" en ese nivel de prioridad.
\end{itemize}

\paragraph{Incorporación del Diagrama de Clases v3.1}

La auditoría de trazabilidad identificó 12 de los 26 RF de esta
matriz (RF-07, RF-08, RF-11, RF-12, RF-14, RF-15, RF-16, RF-17,
RF-18, RF-19, RF-24, RF-26) sin una clase de dominio correspondiente
en el Diagrama de Clases Conceptual, por desactualización del modelo
UML respecto al ERS vigente. Se incorporó el Diagrama de Clases
v3.1 (\texttt{DIAGRAMA DE CLASES REFINADOS v3.1.pdf}), que agrega
cinco clases nuevas --- \textit{Cita}, \textit{Factura},
\textit{SugerenciaIA}, \textit{Consentimiento} y
\textit{SolicitudARCO} --- y resuelve RF-16 reutilizando el método
existente \texttt{HistorialClinico.exportarPDF()}, sin requerir una
clase \textit{Reporte} independiente. Con esta actualización, las 26
filas de la Matriz de Trazabilidad Final quedan con Estado
\textit{Completa} en su columna Clase UML.

\paragraph{Sincronización backlog $\leftrightarrow$ ERS}

Se verificó la sincronización entre el product backlog (tablero
Trello, exportado en \texttt{Trello\_Import\_SGCV\_IA.csv}) y el ERS
vigente: cada uno de los 25 RF activos (RF-01 a RF-24, RF-26; RF-25
reservado) cuenta con al menos una historia de usuario asociada
(HU-01 a HU-25), y cada historia referencia un RF de origen
existente. Sincronización alcanzada: \textbf{100\% (25/25 RF)}, sin
historias huérfanas.

\paragraph{Huérfanos y cadenas rotas}

De los seis hallazgos registrados en la auditoría de trazabilidad,
los dos críticos (colisión de ID, contenido de otro dominio) ya
estaban resueltos al inicio de esta ronda; el hallazgo de clases sin
modelar quedó resuelto con el Diagrama de Clases v3.1 (ver arriba).
No se identificaron requisitos huérfanos (sin fuente) ni artefactos
huérfanos (casos de uso, clases o pruebas sin requisito asociado)
adicionales tras esta corrección.

\subsubsection{Resumen para el informe}

El modelado UML del SGCV-IA mantiene los diez casos de uso de la
iteración anterior (CU-01 a CU-10). El Diagrama de Clases Conceptual
(Tabla 11) se amplió en la versión v3.1 con cinco clases adicionales
--- Cita, Factura, SugerenciaIA, Consentimiento y SolicitudARCO ---,
más el reuso de HistorialClinico.exportarPDF() para RF-16, cerrando
la brecha detectada en la auditoría anterior. Con esta actualización,
los 12 RF que carecían de clase de dominio asociada quedan con
trazabilidad completa hacia el modelo de dominio.

La secuencia CCB (PE4) $\rightarrow$ análisis LOPDP (Usar.tex)
$\rightarrow$ auditoría de trazabilidad (PE5) $\rightarrow$
corrección (\texttt{main.tex}) es, en sí misma, un ejemplo
funcionando de ciclo de gestión de requisitos: un cambio se aprueba
y se documenta (CCB), otro subequipo detecta con evidencia
verificable un hallazgo (auditoría de trazabilidad), y el equipo lo
corrige en la siguiente revisión del documento, dejando además una
nota explícita de por qué se resolvió así. Esto confirma, con un
caso real y completo del propio equipo, que la matriz de
trazabilidad cumple su función cuando se trata como un artefacto
vivo que se audita después de cada línea base, y que el ciclo
completo (detectar, decidir, corregir, verificar) es más valioso que
cualquiera de sus pasos por separado.
